\documentclass[12pt]{article}
\usepackage[margin=0.75in]{geometry}

\newcommand{\duedate}{03/25/2026}
\newcommand{\assignment}{Mem0: Building Production-Ready AI Agents with Scalable Long-Term Memory}

\newcommand{\name}{Aksel Kretsinger-Walters, adk2164}
\newcommand{\email}{adk2164@columbia.edu}

\makeatletter
\def\input@path{{../}{../../}{../../../}}
\makeatother
\input{pset_template_CLMM.tex} %% DO NOT CHANGE THIS LINE

\begin{document}
\psetheader %% DO NOT CHANGE THIS LINE
\section{Summary: Mem0 --- Building Production-Ready AI Agents with Scalable Long-Term Memory}

\subsection{Motivation}
Modern AI agents (e.g., LLM-based assistants) lack persistent, structured long-term memory. While they excel
at short-context reasoning, they struggle to:
\begin{itemize}
		\item Retain user preferences over time
		\item Accumulate knowledge across sessions
		\item Scale memory without degrading performance
\end{itemize}
Mem0 addresses this gap by proposing a \textbf{production-ready memory system} for AI agents that is
scalable, structured, and retrieval-efficient.

\subsection{Core Idea}
Mem0 introduces a modular long-term memory system that:
\begin{itemize}
		\item Stores information as \textbf{discrete memory units}
		\item Continuously \textbf{extracts, updates, and consolidates} memories
		\item Uses \textbf{retrieval mechanisms} to inject relevant memory into prompts
\end{itemize}

The key principle is separating:
\begin{itemize}
		\item \textbf{Working memory} (LLM context window)
		\item \textbf{Long-term memory} (external persistent store)
\end{itemize}

\subsection{Memory Lifecycle}
Mem0 formalizes memory as a pipeline with three stages:

\paragraph{1. Extraction}
Relevant information is extracted from interactions using an LLM. This includes:
\begin{itemize}
		\item User preferences
		\item Facts
		\item Behavioral patterns
\end{itemize}

\paragraph{2. Storage}
Memories are stored in structured form, typically as:
\begin{itemize}
		\item Key-value pairs
		\item Embeddings for semantic retrieval
		\item Metadata (timestamps, importance, etc.)
\end{itemize}

\paragraph{3. Retrieval}
At inference time, relevant memories are retrieved via:
\begin{itemize}
		\item Semantic search (vector similarity)
		\item Filtering based on context or recency
\end{itemize}
These are then injected into the LLM prompt.

\subsection{Memory Operations}
Mem0 emphasizes that memory is \textbf{dynamic}, supporting:
\begin{itemize}
		\item \textbf{Add}: store new information
		\item \textbf{Update}: refine or overwrite existing memories
		\item \textbf{Delete}: remove stale or incorrect memories
\end{itemize}

This avoids unbounded growth and maintains memory quality.

\subsection{Scalability Challenges}
The paper highlights several challenges in scaling memory:
\begin{itemize}
		\item \textbf{Memory explosion}: naive accumulation leads to inefficiency
		\item \textbf{Retrieval noise}: irrelevant memories degrade performance
		\item \textbf{Latency}: large memory stores slow down inference
\end{itemize}

Mem0 mitigates these via:
\begin{itemize}
		\item Memory filtering and ranking
		\item Periodic consolidation
		\item Efficient vector databases
\end{itemize}

\subsection{System Design}
Mem0 is designed as a \textbf{production system}, not just a research prototype:
\begin{itemize}
		\item Modular architecture (plug-and-play memory backend)
		\item Compatible with existing LLM pipelines
		\item Supports real-time updates
\end{itemize}

\subsection{Key Contributions}
\begin{itemize}
		\item A practical framework for long-term memory in AI agents
		\item A full memory lifecycle (extract, store, retrieve, update)
		\item Emphasis on scalability and production constraints
		\item Demonstration of improved personalization and task performance
\end{itemize}

\subsection{Takeaways}
Mem0 shows that:
\begin{itemize}
		\item Long-term memory is essential for production-grade AI agents
		\item Memory must be \textbf{structured, selective, and dynamic}
		\item Retrieval quality is as important as model capability
\end{itemize}

Overall, the paper bridges the gap between LLM capabilities and real-world agent requirements by treating
memory as a first-class system component.
\end{document}
